% Intel Corporation — Software Engineer, Agent Harness — Cover Letter
% Aditya Kulkarni | 2026-06-21
\documentclass[11pt]{article}
\usepackage[left=0.75in,top=0.75in,right=0.75in,bottom=0.75in]{geometry}
\usepackage{fontspec}
\setmainfont{Times New Roman}
\usepackage{fontawesome5}
\usepackage{xcolor}
\usepackage{enumitem}
\usepackage{hyperref}
\pagestyle{empty}

% Intel blue accent
\definecolor{accent}{HTML}{0071C5}

\hypersetup{colorlinks=true, urlcolor=accent}

\begin{document}

%--- Header ---
{\large\textbf{Aditya Kulkarni}} \hfill
{\small
  \textcolor{accent}{\faPhone}\ (281) 876-7066 \quad
  \textcolor{accent}{\faEnvelope}\ adityakulkarnius@gmail.com \quad
  \textcolor{accent}{\faLinkedin}\ \href{https://linkedin.com/in/aditya-kulkarni-355b81217}{linkedin.com/in/aditya-kulkarni-355b81217}
}

{\small
  \textcolor{accent}{\faGithub}\ \href{https://github.com/Aditya-k24}{github.com/Aditya-k24} \quad
  \textcolor{accent}{\faGlobe}\ \href{https://kulkarniaditya.com}{kulkarniaditya.com}
}

\vspace{4pt}
\textcolor{accent}{\rule{\linewidth}{0.5pt}}
\vspace{8pt}

\today

\vspace{8pt}
Hiring Team, Intel Corporation\\
Software Engineer -- Agent Harness (JR0284869)\\
Santa Clara, CA

\vspace{10pt}
Dear Hiring Team,

\vspace{8pt}
I'm applying for the Software Engineer -- Agent Harness role. I've spent the last year building the exact stack Intel is describing: MCP servers for tool schemas and context injection, Temporal.io for durable agent orchestration, and a Kubernetes-native inference platform that routes across multiple model backends.

\vspace{8pt}
\textbf{What I bring to this role:}

\begin{itemize}[leftmargin=*, label={\textcolor{accent}{\faAngleRight}}, itemsep=4pt]
  \item \textbf{Agent harness engineering, end-to-end.} At Isomer AI, I built MCP servers in both Python and TypeScript that power configurable multi-model LLM agent pipelines (Claude, GPT-4o, Gemini) with typed field extraction schemas, context injection, and Playwright browser automation -- deployed to 30+ enterprise customers. I also engineered the Temporal.io workflow engine driving these pipelines: 3-state versioning (Draft\,\textrightarrow\,Live\,\textrightarrow\,Archived), canary traffic splits, QA sampling, and CRITICAL/HIGH/STANDARD risk signal detection.

  \item \textbf{LLM inference infrastructure.} CortexQ is a Kubernetes-native LLM inference platform I built independently: a FastAPI router with a 3-state circuit breaker and EWMA latency-aware load balancing across Claude, GPT-4o, and Gemini backends; a KEDA autoscaler that scales model pods from 0 to N on Redis queue depth; and a Python Kopf CRD operator (LLMDeployment) that reconciles Deployments, Services, and ScaledObjects. This maps directly to Intel's local+cloud hybrid intelligence architecture -- the same scale-to-zero and dynamic routing logic applies to edge inference.

  \item \textbf{Model benchmarking and evaluation.} My LLM Eval Harness compares Claude Sonnet\,4.6, GPT-4o, and Gemini\,2.5 Pro on a 2,000-example frozen test set across grounded QA, multi-hop QA, and claim verification -- tracking ExactMatch, TokenF1, grounded hallucination rate, p50/p95 latency, and cost/call with paired bootstrap CIs and McNemar significance testing. This is the kind of rigorous, vendor-neutral evaluation Intel's agent harness needs.
\end{itemize}

\vspace{8pt}
I'm graduating from NC State's MS CS program in May 2026 and am ready to relocate to Santa Clara immediately. I would welcome the chance to discuss how my background maps to Intel's agent harness work.

\vspace{10pt}
Sincerely,

\vspace{4pt}
\textbf{Aditya Kulkarni}

\end{document}
